%%%%%%%%%%%%%%%%%%%%%%%%%%%%%%%%%%%%%%%%%%%%%%%%%%%%%%%%%%%%%%%%%%%%%%%%%%%%
% FRONT MATTER FOR OECP PAPER
%%%%%%%%%%%%%%%%%%%%%%%%%%%%%%%%%%%%%%%%%%%%%%%%%%%%%%%%%%%%%%%%%%%%%%%%%%%%

\documentclass[12pt]{article}

%---------------------------------------------------------------------------
% Packages
%---------------------------------------------------------------------------
\usepackage[margin=1in]{geometry}
\usepackage{setspace}
\usepackage{amsmath, amssymb, amsfonts}
\usepackage{graphicx}
\usepackage{hyperref}
\usepackage{booktabs}
\usepackage{csquotes}

% Bibliography setup
\usepackage[backend=bibtex,style=authoryear,sorting=nyt,maxcitenames=2]{biblatex}
\addbibresource{oecp_refs.bib}

\onehalfspacing

%---------------------------------------------------------------------------
% Title Information
%---------------------------------------------------------------------------

\title{\Large\bfseries
Open Economy Carbon Pricing:
A Trade-Based Framework for Efficient, Politically Feasible,
and Rapid Decarbonisation}

\author{
    Stephen J.\ Stretton\thanks{Climate Economist, World Bank Group.  
    All views expressed are strictly personal and do not represent the World Bank, 
    its Executive Directors, or the countries they represent.}
}

\date{\today}

%---------------------------------------------------------------------------
% Begin Document
%---------------------------------------------------------------------------

\begin{document}

\maketitle
\thispagestyle{empty}

% ==================== ABSTRACT ====================
\begin{abstract}
Carbon pricing remains too weak in most countries to deliver rapid decarbonisation. We argue this weakness is structural, arising from three constraints: (1) the trade-exposure constraint, (2) the political-economy constraint, and (3) the primary-to-quaternary structure of energy transition. Conventional carbon pricing (CCP) applies a source-based tax to domestic production, creating competitiveness risks and inviting a race to the bottom similar to corporate income taxation. Open Economy Carbon Pricing (OECP) restructures carbon pricing into a destination-based framework, pairing a production-side feebate with a consumption-based carbon charge that applies symmetrically to imports and domestic goods while zero-rating exports. This enables a race to the top akin to the rise of VAT. We present the simplest integrated implementation mechanism using CBAM logic and common external tariff structures. Appendices provide a simple trade model and a simple political-feasibility model.
\end{abstract}

\newpage

\tableofcontents

%%%%%%%%%%%%%%%%%%%%%%%%%%%%%%%%%%%%%%%%%%%%%%%%%%%%%%%%%%%%%%%%%%%%%%%%%%%%
% PART I — INTRODUCTION, LITERATURE, STRUCTURAL CONSTRAINTS, OECP FRAMEWORK
% AND THE TWO-PLUS-ONE FUNCTIONS OF CARBON PRICING
%%%%%%%%%%%%%%%%%%%%%%%%%%%%%%%%%%%%%%%%%%%%%%%%%%%%%%%%%%%%%%%%%%%%%%%%%%%%
\newpage
\section{Introduction}

Carbon pricing is widely regarded as a foundational instrument for climate policy. 
Yet after three decades of implementation, effective carbon prices remain far below 
levels required for accelerated decarbonisation. The common explanation is political 
resistance, but this paper argues that the true problem is structural: 
\textit{conventional carbon pricing (CCP) is a source-based production tax}, 
and this architecture embeds three deep constraints that no amount of political 
effort can fully overcome.

First, CCP induces a persistent \textit{trade constraint}. 
Because the tax is applied where production occurs, it raises domestic production 
costs, weakens export competitiveness, and encourages import substitution 
from jurisdictions with lower carbon constraints. This creates an international 
race to the bottom, analogous to the erosion of corporate income tax rates.

Second, CCP suffers from a severe \textit{political-economy constraint}. 
Carbon taxes and ETS compliance costs are highly salient for both producers and 
households: producers face higher average costs and lobby against increases, 
while households face higher retail prices and vote accordingly.
The politically feasible carbon price is therefore limited well below 
the level required to induce capital turnover in primary sectors.

Third, CCP creates a \textit{transition constraint}. 
Primary sectors such as steel, cement, electricity, refining, and heavy transport 
are capital-intensive and slow to respond to marginal price changes.
A carbon price high enough to accelerate early switching 
is politically infeasible under CCP’s architecture.

This paper proposes a structural redesign: 
\textbf{Open Economy Carbon Pricing (OECP)},  
which resolves all three constraints by separating the 
functions of carbon pricing across two instruments:
\begin{enumerate}
    \item a \textbf{revenue-neutral production feebate} that strengthens 
    clean–dirty competition without raising average production costs; and
    \item a \textbf{destination-based consumption carbon charge} applied to all 
    goods consumed domestically, with exports zero-rated.
\end{enumerate}

This architecture mirrors the transformation that the Value Added Tax (VAT) 
achieved when replacing source-based turnover taxes: shifting taxation to the 
destination basis eliminates competitiveness harms and enables a ``race to the top.''  

Before developing the detailed implementation (Part II) and models (Part III),  
Part I sets out four foundational components:
\begin{enumerate}
    \item the relevant literatures;  
    \item the structural constraints of CCP;  
    \item the OECP framework; and  
    \item the \textbf{two-plus-one functions of carbon pricing}, 
          a new conceptual lens that clarifies why CCP fails and why OECP works.
\end{enumerate}

%%%%%%%%%%%%%%%%%%%%%%%%%%%%%%%%%%%%%%%%%%%%%%%%%%%%%%%%%%%%%%%%%%%%%%%%%%%%
\section{Literature Review}

Open Economy Carbon Pricing (OECP) sits at the intersection of climate economics,
international trade, environmental taxation, and public finance.  
The relevant literature spans six strands:
(1) competitiveness and leakage under unilateral carbon pricing;  
(2) benchmarked instruments, feebates, and output-based rebating;  
(3) border carbon adjustments and carbon-adjusted tariff structures;  
(4) consumption-based greenhouse gas accounting;  
(5) consumption-based environmental taxation;  
and (6) political acceptability and destination-based taxation in public finance.  

This section reviews each in turn and highlights how OECP synthesises their insights.

%%%%%%%%%%%%%%%%%%%%%%%%%%%%%%%%%%%%%%%%%%%%%%%%%%%%%%%%%%%%%%%%%%%%%%%%%%%%
\subsection{Competitiveness, Leakage, and Unilateral Carbon Pricing}

A substantial body of work has shown that conventional, source-based carbon pricing
creates competitiveness concerns and carbon leakage when implemented unilaterally.
Seminal contributions include \textcite{FischerFox2007, FischerFox2012},  
who show analytically and numerically that carbon pricing raises costs in  
energy-intensive, trade-exposed (EITE) sectors and shifts production abroad.  
General equilibrium analyses such as \textcite{BoehringerRutherford2016}  
confirm that leakage rates can be substantial even under moderate carbon prices.

Theoretical research in environmental trade theory  
(\textcite{Hoel1996}; \textcite{CopelandTaylor2005})  
demonstrates that open economies internalise only a fraction of the environmental  
externality, setting carbon prices below global optima.  
\textcite{AldyPizer2015} further show that carbon-price differentials across  
countries can persist even under cooperative arrangements when carbon pricing  
remains production-based.

This literature establishes that **source-based carbon pricing induces a structural 
race to the bottom**, creating the core policy problem that OECP seeks to resolve.

%%%%%%%%%%%%%%%%%%%%%%%%%%%%%%%%%%%%%%%%%%%%%%%%%%%%%%%%%%%%%%%%%%%%%%%%%%%%
\subsection{Feebates, Benchmarks, and Output-Based Rebating}

Benchmark-based policy instruments offer relative incentives without raising average
costs.  
\textcite{FischerFox2012} and related work (\textcite{Fischer2019})  
show that output-based rebating (OBR) protects competitiveness while  
preserving some marginal abatement incentive.  
Feebate systems—analysed in \textcite{BoehringerLange2005} and  
applied in multiple sectors including vehicles, electricity, and industry  
(\textcite{HollandMansurYates2020})—create linear transfers based on  
emissions intensity relative to a benchmark.

The key insight is that feebates can provide robust **relative** incentives to
decarbonise while avoiding **absolute** cost increases, thereby eliminating the 
primary source of industrial opposition to carbon pricing.

OECP generalises this approach by using a sector-wide revenue-neutral feebate
as the production-side instrument, enabling high incentive rates without harming 
competitiveness.

%%%%%%%%%%%%%%%%%%%%%%%%%%%%%%%%%%%%%%%%%%%%%%%%%%%%%%%%%%%%%%%%%%%%%%%%%%%%
\subsection{Border Carbon Adjustments and Carbon-Adjusted Tariffs}

Border carbon adjustments (BCAs) have been widely proposed as a response to
competitiveness and leakage concerns.  
Foundational work by \textcite{IsmerNeuhoff2007} and  
\textcite{MattooSubramanian2009} examines how import charges and export rebates 
could mimic the effects of a domestic carbon price at the border.  
Subsequent work synthesising legal and economic perspectives  
(\textcite{Mehling2019}; \textcite{DroegeFischer2020})  
identifies persistent challenges:  
verifying embodied emissions, assessing foreign climate policies,  
treating free allocation, and ensuring WTO consistency.

More recent literature extends border adjustments toward **environmental tariff
reforms**, often referred to as “green tariffs.”  
This work includes environmental border charges  
(\textcite{LockwoodWhalley2008}), early proposals for  
climate-related tariff adjustments (\textcite{Low2008}; \textcite{Pauwlyn2007}),  
and analyses of integrating carbon criteria into the EU’s  
Common External Tariff (\textcite{Colares2014}).  
Studies on differentiated tariffs based on verified low-carbon production  
(\textcite{ElectriciteDeFrance2019}) and HS-level embodied-emissions adjustments  
(\textcite{MoranWood2020}; \textcite{Brandi2023})  
demonstrate that customs-based climate policy is legally and operationally feasible.

These contributions highlight two insights reflected in OECP:
(1) customs systems can incorporate carbon content using benchmark intensities;  
(2) destination-based adjustments can be more transparent and administratively 
simpler than CBAM-like schemes that track foreign policies.

%%%%%%%%%%%%%%%%%%%%%%%%%%%%%%%%%%%%%%%%%%%%%%%%%%%%%%%%%%%%%%%%%%%%%%%%%%%%
\subsection{Consumption-Based Emissions Accounting}

A substantial literature in industrial ecology and climate policy argues that
emissions responsibility should follow consumption.  
\textcite{Peters2008}, \textcite{PetersHertwich2008},  
and \textcite{HertwichPeters2009} demonstrate that many high-income countries have  
consumption-based emissions far exceeding territorial emissions, due to  
outsourcing of carbon-intensive production.  
\textcite{JakobSteckel2014} and \textcite{JakobEdenhofer2014} emphasise that  
consumption-based frameworks better reflect equity considerations  
and the structure of global value chains.

This literature motivates OECP’s destination-basis approach:  
carbon charges should apply where goods are consumed, not where they are produced.

%%%%%%%%%%%%%%%%%%%%%%%%%%%%%%%%%%%%%%%%%%%%%%%%%%%%%%%%%%%%%%%%%%%%%%%%%%%%
\subsection{Consumption-Based Environmental Taxation and Carbon-Adjusted Tariff Systems}

Consumption-based environmental taxation provides real-world examples of how
environmental externalities can be incorporated into prices downstream.

Empirical evidence comes from:
\begin{itemize}
    \item Nordic NO$_x$ charges rebated according to performance benchmarks  
    (\textcite{BruvollLarsen2004});
    \item landfill taxes and waste-disposal levies in the UK  
    (\textcite{Elliott2011});
    \item packaging, pesticide, and chemical taxes across Europe  
    (\textcite{OECD2016});
    \item extended producer-responsibility schemes such as WEEE and battery directives  
    (\textcite{Tojo2004});
    \item utility and retail energy levies including supplier-obligation schemes  
    (\textcite{Hood2011}).
\end{itemize}

Recent research on consumption-based carbon pricing examines:
retail electricity carbon charges (\textcite{VogtSchilb2022});  
point-of-sale carbon footprint fees (\textcite{Summerfield2019});  
and platform-based carbon charges.

Parallel work considers **carbon-adjusted tariffs**,  
a consumption-based mechanism implemented at the border:  
environmental border charges \textcite{LockwoodWhalley2008};  
tariff refinement based on climate criteria \textcite{Colares2014};  
verified low-carbon tariff differentiation \textcite{ElectriciteDeFrance2019};  
and embodied-emissions adjustments to customs schedules  
\textcite{MoranWood2020}; \textcite{Brandi2023}.

These literatures collectively show that consumption-based carbon taxation is
administratively feasible, compatible with trade rules, and increasingly embedded in 
real-world environmental policy.

OECP builds directly on this foundation by combining a destination-based  
consumption charge with benchmark-based tariff adjustments.

%%%%%%%%%%%%%%%%%%%%%%%%%%%%%%%%%%%%%%%%%%%%%%%%%%%%%%%%%%%%%%%%%%%%%%%%%%%%
\subsection{Political Acceptability, Revenue Use, and Public Finance Analogies}

Political-economy research shows that public support for carbon pricing depends
heavily on how revenues are used.  
\textcite{Metcalf2009, Metcalf2017} emphasise transparent, equitable recycling.  
\textcite{ParryWilliams1999} analyse efficiency–equity trade-offs in revenue 
use, while \textcite{Carbone2013} demonstrate the distributional impacts of  
alternative recycling mechanisms.  
Survey evidence synthesised in \textcite{Klenert2018} shows that visible household 
transfers substantially increase acceptability.

The public-finance literature provides an important analogy:  
destination-based taxes are less distortionary than source-based taxes.  
Empirical studies on VAT adoption (\textcite{KeenLockwood2006, KeenLockwood2010})  
show that VAT improved revenue stability and encouraged upward convergence  
because it taxes consumption rather than production.  
Theoretical work on destination-based corporate taxation 
(\textcite{Auerbach2010}) shows similar advantages.

OECP applies these insights to carbon:  
it separates revenue generation (via the consumption charge)  
from incentive provision (via feebates),  
mirroring the advantages of moving from source-based to destination-based taxation.

%%%%%%%%%%%%%%%%%%%%%%%%%%%%%%%%%%%%%%%%%%%%%%%%%%%%%%%%%%%%%%%%%%%%%%%%%%%%
% END LITERATURE REVIEW
%%%%%%%%%%%%%%%%%%%%%%%%%%%%%%%%%%%%%%%%%%%%%%%%%%%%%%%%%%%%%%%%%%%%%%%%%%%%

%%%%%%%%%%%%%%%%%%%%%%%%%%%%%%%%%%%%%%%%%%%%%%%%%%%%%%%%%%%%%%%%%%%%%%%%%%%%
\section{The Two-Plus-One Functions of Carbon Pricing}

A major contribution of this paper is clarifying that carbon pricing 
has \textbf{two established functions} and \textbf{one missing function} that 
current policy does not address.

\subsection{Function 1: Transition (Economic Efficiency)}

The first purpose of carbon pricing is to induce a shift from high-emission 
to low-emission technologies.  
This requires strong \textit{relative} price signals, not necessarily high 
\textit{average} cost increases.  
Feebates are ideal for this function.

\subsection{Function 2: Polluter-Pays Responsibility}

The second function is normative: 
entities that cause climate damage should pay for that damage.  
Under CCP, this function is combined with the transition function within a 
single instrument, creating tension between high prices for justice and 
low prices for feasibility.  
OECP separates these functions: the consumption charge handles responsibility, 
while the feebate handles incentives.

\subsection{Function 3 (Missing): Remediation or Compensation}

In moral philosophy and tort law, responsibility must be paired with:
\begin{itemize}
    \item remediation of harm; or  
    \item compensation to harmed parties.
\end{itemize}
Current systems do neither adequately.  
OECP, by relocating revenue generation to consumption, frees sufficient revenue 
for a future system of remediation and compensation.  
Designing such a system is beyond this paper’s scope but constitutes a 
natural extension of the OECP framework.

\subsection{Why Combining Functions Breaks CCP}

Under CCP:
\begin{itemize}
    \item Transition requires high prices;  
    \item Polluter-pays justice also requires high prices;  
    \item Political feasibility constrains prices to low levels.
\end{itemize}
Combining functions within a single tax collapses under political pressure.  
OECP succeeds because it separates functions.


%%%%%%%%%%%%%%%%%%%%%%%%%%%%%%%%%%%%%%%%%%%%%%%%%%%%%%%%%%%%%%%%%%%%%%%%%%%%
\section{The Structural Constraints of Conventional Carbon Pricing}

\subsection{The Trade Constraint}

Source-based carbon pricing raises both marginal and average production costs 
for domestic firms.  
Exports become less competitive, and imports—often from more carbon-intensive 
producers—gain an advantage.  
Even with CBAM, domestic producers remain at a disadvantage in export markets.  
This creates structural leakage and a persistent race to the bottom in carbon prices.

\subsection{The Political-Economy Constraint}

Political feasibility under CCP requires keeping:
\begin{itemize}
    \item industrial opposition manageable; and
    \item household price impacts tolerable.
\end{itemize}
A uniform carbon tax or ETS increases both.  
Empirically, feasible prices remain in the €30–€60/tCO$_2$ range, well below 
the level required for rapid structural transition.

\subsection{The Transition Constraint}

Primary sectors require capital turnover, not marginal adjustment.  
The carbon price required to shift investment in steelmaking, cement production, 
refining or power generation is far higher than the price politically feasible 
under CCP.  
Thus, CCP is structurally mismatched to the energy system’s 
primary–tertiary structure.


%%%%%%%%%%%%%%%%%%%%%%%%%%%%%%%%%%%%%%%%%%%%%%%%%%%%%%%%%%%%%%%%%%%%%%%%%%%%
\section{Interoducing Open-Economy Carbon Pricing (OECP) Framework}

OECP replaces Conventional Carbon Pricing (CCP) -- which has a single, overloaded instrument  -- with a two-instrument system separating incentives (production feebate) from damage remediation (consumption based carbon pricing).

\subsection{Production Feebate}

Each sector sets a benchmark emissions intensity $\bar e$.  
A firm with emissions $e_i$ receives a subsidy or pays a fee:
\[
F_i = t(e_i - \bar e).
\]
This preserves relative incentives while keeping sector-average costs unchanged.  
This eliminates export disadvantages and encourages technological switching.

\subsection{Destination-Based Consumption Charge}

All goods consumed domestically face a charge $t \times m_j$, 
where $m_j$ is embodied emissions.  
Imports and domestic goods are treated symmetrically; exports are zero-rated.  
This mirrors VAT and restores trade neutrality.

\subsection{Instrument Interaction}

The feebate strengthens domestic decarbonisation competition.  
The consumption charge:
\begin{itemize}
    \item raises revenue,  
    \item applies polluter-pays principles,  
    \item and eliminates trade distortions.
\end{itemize}
Together, they raise feasible incentive rates and avoid leakage.

\subsection{Forward Links}

Appendix A formalises the trade model underlying this logic.  
Appendix B develops the political feasibility model.  
Appendix C provides the full implementation details.

%%%%%%%%%%%%%%%%%%%%%%%%%%%%%%%%%%%%%%%%%%%%%%%%%%%%%%%%%%%%%%%%%%%%%%%%%%%%
\section{Accelerated Transition: The Primary--Tertiary Argument}

Carbon pricing must align with the structure of the economy:
\begin{itemize}
    \item primary sectors (fuels, steel, cement),
    \item secondary sectors (manufactured goods),
    \item tertiary sectors (services).
\end{itemize}

Most emissions originate upstream in primary sectors, yet downstream consumers drive
demand.  
Conventional carbon pricing, by taxing only domestic production, fails to accelerate
capital turnover in primary sectors because it:
\begin{itemize}
    \item raises average costs (thus generating political resistance),
    \item but fails to provide strong relative incentives,
    \item and cannot reach required levels because of trade and voter constraints.
\end{itemize}

OECP resolves this:
\begin{itemize}
    \item strong feebate incentives accelerate switching,
    \item destination-basis charges allow high effective carbon prices,
    \item export neutrality prevents industrial pushback,
    \item consumption charges align downstream demand with upstream transition.
\end{itemize}

This restores coherence between climate policy design and energy-system structure.


\section{Implementation of OECP}

OECP consists of three administratively simple elements:
\begin{enumerate}
    \item a revenue-neutral production feebate,
    \item a destination-based consumption carbon charge,
    \item and border adjustments for imports and exports.
\end{enumerate}
This section shows how these components can be implemented using existing
institutional infrastructure, building on but simplifying the logic of the EU CBAM
and universal VAT systems.

The overall objective is to ensure:
\begin{itemize}
    \item strong decarbonisation incentives for producers,
    \item equal treatment of domestic and foreign goods,
    \item preserved export competitiveness,
    \item and politically sustainable revenue generation.
\end{itemize}

Implementation details requiring deeper technical specification are deferred to
Appendix C (Tax Implementation), but summarised here for clarity.

%%%%%%%%%%%%%%%%%%%%%%%%%%%%%%%%%%%%%%%%%%%%%%%%%%%%%%%%%%%%%%%%%%%%%%%%%%%%
\subsection{Production Feebate: Incentives Without Cost Increases}

In each sector, a benchmark emissions intensity $\bar e$ is chosen.  
Any producer $i$ with emissions $e_i$ receives a feebate:
\[
F_i = t (e_i - \bar e).
\]

The feebate has four critical properties:
\begin{enumerate}
    \item It is approximately revenue-neutral at the sector level.
    \item It creates strong relative incentives to decarbonise.
    \item It does not raise average production costs.
    \item It leaves export prices unchanged.
\end{enumerate}

This is the core instrument that resolves CCP’s competitiveness constraint.
Because average costs do not increase, firms do not lobby against high incentive rates.
Producers remain competitive in global markets even when $t$ is high.

%%%%%%%%%%%%%%%%%%%%%%%%%%%%%%%%%%%%%%%%%%%%%%%%%%%%%%%%%%%%%%%%%%%%%%%%%%%%
\subsection{Consumption Carbon Charge: Revenue and Responsibility}

All goods consumed domestically—regardless of origin—bear a carbon charge:
\[
C_j = t \cdot m_j,
\]
where $m_j$ is embodied emissions in product $j$.  

This charge:
\begin{itemize}
    \item applies identically to imports and domestic goods,
    \item ensures polluter-pays responsibility,
    \item raises revenue without imposing costs on exports,
    \item is compatible with WTO norms under GATT Article II:2(a),
    \item and mirrors VAT systems, which tax consumption rather than production.
\end{itemize}

The consumption charge is the locus of revenue generation within OECP.
It replaces the revenue-raising function of CCP and allows future work
(Section~\ref{sec:remediationfuture}) to allocate revenue for remediation and compensation.

%%%%%%%%%%%%%%%%%%%%%%%%%%%%%%%%%%%%%%%%%%%%%%%%%%%%%%%%%%%%%%%%%%%%%%%%%%%%
\subsection{Import Adjustments: A Simplified CBAM}

Imports are charged based on embodied emissions.  
The system is intentionally simple:

\begin{enumerate}
    \item Importer declares either verified emissions or accepts a conservative default.
    \item Customs applies $t \cdot m_j$, using:
    \begin{itemize}
        \item a reference intensity (e.g.\ China or global high-bound),
        \item a country-specific adjustment,
        \item optional producer-specific data for large firms.
    \end{itemize}
\end{enumerate}

This structure avoids CBAM’s complexity:
\begin{itemize}
    \item no tracking of foreign free allocations,
    \item no foreign ETS equivalence calculations,
    \item and no need for the EU-style “carbon certificate market.”
\end{itemize}

Implementation can follow:
\begin{itemize}
    \item \textbf{category-level intensities} for EITE sectors (steel, cement, aluminium),
    \item \textbf{value-based intensities} for general goods,
    \item \textbf{physical-unit intensities} for bulk commodities (see Appendix C).
\end{itemize}

%%%%%%%%%%%%%%%%%%%%%%%%%%%%%%%%%%%%%%%%%%%%%%%%%%%%%%%%%%%%%%%%%%%%%%%%%%%%
\subsection{Export Treatment: VAT-Style Zero Rating}

Exports face \emph{no} consumption charge.  
Producers still participate in the feebate but pay no absolute carbon tax on exports.

Thus exports remain competitive, and domestic producers are not harmed by high
carbon incentive rates.  
This preserves neutrality across jurisdictions and ensures WTO compatibility.

%%%%%%%%%%%%%%%%%%%%%%%%%%%%%%%%%%%%%%%%%%%%%%%%%%%%%%%%%%%%%%%%%%%%%%%%%%%%
\subsection{Common External Tariff Variant}

In customs unions, a \textbf{Carbon-Adjusted Common External Tariff (CACET)} can
simplify implementation:

\begin{itemize}
    \item A base tariff applies to all countries.
    \item A carbon adjustment is added:
    \[
    \text{CAT}_j = t \cdot (m_j - m_{\text{benchmark}}).
    \]
    \item Countries with verified low-carbon production receive reduced charges.
\end{itemize}

This approach is particularly suitable for the EU, ASEAN, EAC, and MERCOSUR.

%%%%%%%%%%%%%%%%%%%%%%%%%%%%%%%%%%%%%%%%%%%%%%%%%%%%%%%%%%%%%%%%%%%%%%%%%%%%
\subsection{International Incentives: Why OECP Encourages a Race to the Top}

OECP changes the international incentive structure by rewarding lower-carbon
production regardless of jurisdiction.

\paragraph{Under CCP:}
A country that decarbonises its industry faces higher domestic costs and no
advantage for exports.  
Thus the incentive is weak.

\paragraph{Under OECP:}
A country that decarbonises:
\begin{itemize}
    \item receives lower import charges in foreign markets (as foreign consumers pay less),
    \item secures competitive gains for its domestic producers,
    \item and retains strong feebate incentives internally.
\end{itemize}

This restores the same logic that made VAT universally adoptable: destination-basis
taxes remove competitiveness harms and induce upward convergence.

%%%%%%%%%%%%%%%%%%%%%%%%%%%%%%%%%%%%%%%%%%%%%%%%%%%%%%%%%%%%%%%%%%%%%%%%%%%%
\subsection{Transfers and Feebate Gaps}

A subtle issue arises when the feebate incentive becomes strong enough to trigger
switching to a technology that is initially more expensive.  
In this case:
\begin{itemize}
    \item below the switching threshold, feebates generate surplus revenue;
    \item above the threshold, feebates may become revenue-negative.
\end{itemize}

The gap can be financed by:
\begin{itemize}
    \item revenue from the consumption charge, or
    \item international transfers (if used in developing countries),
    \item multilateral development bank facilities,
    \item climate finance arrangements.
\end{itemize}

In a simple two-technology model:
\begin{itemize}
    \item if clean technology is 10\% more expensive,
    \item and if the effective feebate rate reduces average cost increases to 2--3\%,
\end{itemize}
then the transfer requirement is small relative to the total revenue potential of
the consumption charge.

Appendix A provides a formal derivation of this result.


%%%%%%%%%%%%%%%%%%%%%%%%%%%%%%%%%%%%%%%%%%%%%%%%%%%%%%%%%%%%%%%%%%%%%%%%%%%%
\section{Link to Appendices}

Part III provides:
\begin{itemize}
    \item Appendix~A: A formalised economic model of trade and transition.
    \item Appendix~B: A detailed political-economy model showing the feasibility
    frontier and how OECP expands it.
    \item Appendix~C: A complete ``Tax Implementation'' appendix covering import
    charges, country-specific adjustment, product granularity, and export zero-rating.
\end{itemize}

%%%%%%%%%%%%%%%%%%%%%%%%%%%%%%%%%%%%%%%%%%%%%%%%%%%%%%%%%%%%%%%%%%%%%%%%%%%%
% PART III — APPENDICES (ECONOMIC MODEL, POLITICAL MODEL, TAX IMPLEMENTATION)
%%%%%%%%%%%%%%%%%%%%%%%%%%%%%%%%%%%%%%%%%%%%%%%%%%%%%%%%%%%%%%%%%%%%%%%%%%%%

\appendix

%%%%%%%%%%%%%%%%%%%%%%%%%%%%%%%%%%%%%%%%%%%%%%%%%%%%%%%%%%%%%%%%%%%%%%%%%%%%
\section{Appendix A: Economic Model of Trade, Transition, and OECP}

This appendix presents a minimal open-economy model showing:

\begin{enumerate}
    \item why conventional carbon pricing (CCP) creates trade distortions,
    \item why feebates preserve competitiveness,
    \item how destination-based charges eliminate leakage,
    \item and how OECP raises feasible incentive rates.
\end{enumerate}

The model favours simplicity over formal complexity and is intended to 
provide clear intuition for the results used in the main text.

%%%%%%%%%%%%%%%%%%%%%%%%%%%%%%%%%%%%%%%%%%%%%%%%%%%%%%%%%%%%%%%%%%%%%%%%%%%%
\subsection{A.1 Basic Setup}

Two countries: Home (H) and Foreign (F).  
One homogeneous good is produced using one of two technologies:

\begin{itemize}
    \item \emph{Dirty technology (D)}: cost $c_D$, emissions $e_D$.
    \item \emph{Clean technology (C)}: cost $c_C > c_D$, emissions $e_C < e_D$.
\end{itemize}

Demand in Home is $Q_H(p)$; demand in Foreign is $Q_F(p)$.

Foreign imposes no carbon constraints unless otherwise stated.

Trade is frictionless, and firms are price-takers.

%%%%%%%%%%%%%%%%%%%%%%%%%%%%%%%%%%%%%%%%%%%%%%%%%%%%%%%%%%%%%%%%%%%%%%%%%%%%
\subsection{A.2 Conventional Carbon Pricing}

Under a source-based carbon tax $t$, Home producers face cost:

\[
p_i^{\text{CCP}} = c_i + t e_i.
\]

Thus:

\begin{itemize}
    \item Clean production becomes more competitive relative to dirty production.
    \item \textbf{But} both clean and dirty production become less competitive 
          relative to Foreign producers.
\end{itemize}

Exports shrink; imports expand.  
Leakage occurs if the increase in Foreign output exceeds the decline in Home output.

The equilibrium carbon price under CCP is limited by:

\[
t \le t^{\max} = \min(\tau_I, \tau_V),
\]

where $\tau_I$ is the industry tolerance threshold, and $\tau_V$ the voter threshold.

%%%%%%%%%%%%%%%%%%%%%%%%%%%%%%%%%%%%%%%%%%%%%%%%%%%%%%%%%%%%%%%%%%%%%%%%%%%%
\subsection{A.3 OECP Feebate}

Under OECP, producers face a feebate:

\[
F_i = t (e_i - \bar e),
\]

yielding a net cost:

\[
p_i^{\text{FB}} = c_i + t(e_i - \bar e).
\]

If the benchmark $\bar e$ is chosen such that:

\[
\sum_i F_i = 0,
\]

then average production costs \emph{do not change}.  
Thus:

\begin{itemize}
    \item Home producers do not suffer export losses.
    \item Foreign producers gain no competitive advantage.
    \item Strong relative incentives remain for clean producers.
\end{itemize}

%%%%%%%%%%%%%%%%%%%%%%%%%%%%%%%%%%%%%%%%%%%%%%%%%%%%%%%%%%%%%%%%%%%%%%%%%%%%
\subsection{A.4 OECP Consumption Charge}

Home applies a consumption-based carbon charge:

\[
C_j = t \cdot m_j,
\]

to all \emph{goods consumed domestically}, where $m_j$ is embodied emissions.

Thus:

\begin{itemize}
    \item Imports pay $t \cdot m_j$ at the border.
    \item Domestic goods pay $t \cdot m_j$ at point of sale.
    \item Exports pay no charge (zero-rated).
\end{itemize}

%%%%%%%%%%%%%%%%%%%%%%%%%%%%%%%%%%%%%%%%%%%%%%%%%%%%%%%%%%%%%%%%%%%%%%%%%%%%
\subsection{A.5 Results}

\paragraph{Result 1: No leakage}

Because all consumption faces the same charge:

\[
p^{\text{dom}}(j) = p^{\text{imp}}(j),
\]

imports do not displace domestic production.

\paragraph{Result 2: Preservation of export competitiveness}

Because feebates do not raise average costs:

\[
p^{\text{export}}_i = c_i + t (e_i - \bar e) \approx c_i.
\]

Exports remain unaffected.

\paragraph{Result 3: Higher feasible incentive rate}

Since neither households nor industry face large average-cost increases:

\[
t_{\text{OECP}} > t_{\text{CCP}}.
\]

\paragraph{Result 4: Accelerated switching}

The switching condition under OECP becomes:

\[
c_C - c_D < t (e_D - e_C),
\]

which can be satisfied at moderate values of $t$, because producers face 
high relative incentives but no average-cost increases.

%%%%%%%%%%%%%%%%%%%%%%%%%%%%%%%%%%%%%%%%%%%%%%%%%%%%%%%%%%%%%%%%%%%%%%%%%%%%
\section{Appendix B: Political Economy Model}

This appendix develops a simple but more detailed model of political feasibility, 
using two constraints:

\begin{enumerate}
    \item industry tolerance ($\tau_I$),  
    \item voter tolerance ($\tau_V$).
\end{enumerate}

CCP faces both constraints; OECP reshapes them.

%%%%%%%%%%%%%%%%%%%%%%%%%%%%%%%%%%%%%%%%%%%%%%%%%%%%%%%%%%%%%%%%%%%%%%%%%%%%
\subsection{B.1 Utility and Cost Functions}

Let $t$ be the carbon incentive rate.

For industry, the relevant cost under CCP is:

\[
\Delta c_I(t) = t \cdot \bar e_{\text{industry}},
\]

which increases average costs and reduces profits:

\[
\pi_I(t) = \pi_0 - Q \cdot \Delta c_I(t).
\]

Industry accepts CCP if:

\[
\Delta c_I(t) \le \tau_I.
\]

For households, the relevant cost under CCP is:

\[
\Delta c_V(t) = \alpha t,
\]

where $\alpha$ reflects the share of consumption affected and 
salience of energy price changes.

Households accept CCP if:

\[
\Delta c_V(t) \le \tau_V.
\]

Thus:

\[
t_{\text{CCP}}^{\max} = \min(\tau_I/\bar e_{\text{industry}}, \ \tau_V/\alpha).
\]

Empirically, both are low.

%%%%%%%%%%%%%%%%%%%%%%%%%%%%%%%%%%%%%%%%%%%%%%%%%%%%%%%%%%%%%%%%%%%%%%%%%%%%
\subsection{B.2 How OECP Changes the Constraints}

\paragraph{Industry:}

Under OECP:

\[
\Delta c_I^{\text{OECP}}(t) \approx 0.
\]

The feebate preserves average costs:

\[
p_i^{\text{FB}} \approx c_i,
\]

so:

\[
\tau_I^{\text{OECP}} \to \infty \quad \text{(no meaningful constraint)}.
\]

\paragraph{Households:}

The consumption charge raises prices for some goods, but revenue can be 
recycled back to households through:

\begin{itemize}
    \item dividends,
    \item reduced labour taxes,
    \item rebates to low-income households,
    \item support for energy efficiency.
\end{itemize}

Thus the net household cost is:

\[
\Delta c_V^{\text{OECP}}(t) = \alpha t - R(t),
\]

where $R(t)$ is the per-household rebate.

If $R(t)$ is designed to grow with $t$, then:

\[
\Delta c_V^{\text{OECP}}(t) \approx 0.
\]

Thus:

\[
t_{\text{OECP}}^{\max} \gg t_{\text{CCP}}^{\max}.
\]

%%%%%%%%%%%%%%%%%%%%%%%%%%%%%%%%%%%%%%%%%%%%%%%%%%%%%%%%%%%%%%%%%%%%%%%%%%%%
\subsection{B.3 Feasible Policy Region}

The feasible region in $(t, \text{rebate})$ space expands dramatically:

\[
\text{Feasible}_{\text{OECP}} =
\{(t, R): \ \Delta c_I^{\text{OECP}}(t) \approx 0, \ 
          \Delta c_V^{\text{OECP}}(t) \le \tau_V \}.
\]

In graphical terms:
\begin{itemize}
    \item The industry constraint disappears.
    \item The voter constraint becomes a rising function of $R$.
\end{itemize}

Thus OECP opens political space for high incentive rates combined with 
household protection.

%%%%%%%%%%%%%%%%%%%%%%%%%%%%%%%%%%%%%%%%%%%%%%%%%%%%%%%%%%%%%%%%%%%%%%%%%%%%
\section{Appendix C: Tax Implementation}

This appendix provides a full specification of the tax implementation system 
summarised in Part II.

%%%%%%%%%%%%%%%%%%%%%%%%%%%%%%%%%%%%%%%%%%%%%%%%%%%%%%%%%%%%%%%%%%%%%%%%%%%%
\subsection{C.1 Overview of the System}

Imports, domestic production, and exports are treated as follows:

\begin{enumerate}
    \item \textbf{Production feebate}: revenue-neutral, applied to all domestic producers.
    \item \textbf{Consumption charge}: applied to all goods consumed domestically.
    \item \textbf{Import adjustments}: based on embodied emissions.
    \item \textbf{Export zero-rating}: no consumption charge on exports.
\end{enumerate}

%%%%%%%%%%%%%%%%%%%%%%%%%%%%%%%%%%%%%%%%%%%%%%%%%%%%%%%%%%%%%%%%%%%%%%%%%%%%
\subsection{C.2 Country-Level Adjustment}

Each exporting country $k$ receives a default intensity:

\[
m_{jk}^{\text{default}} = m_j^{\text{ref}} + \delta_k,
\]

where:

\begin{itemize}
    \item $m_j^{\text{ref}}$ is a high-bound reference intensity (e.g.\ China),
    \item $\delta_k$ is an adjustment for country $k$'s sectoral emissions profile.
\end{itemize}

Large firms can optionally submit producer-specific data $m_{jk}^{\text{firm}}$.

%%%%%%%%%%%%%%%%%%%%%%%%%%%%%%%%%%%%%%%%%%%%%%%%%%%%%%%%%%%%%%%%%%%%%%%%%%%%
\subsection{C.3 Product Granularity Options}

Three methods:

\paragraph{(1) Category-level}:  
HS-code based categories with default intensities.

\paragraph{(2) Physical-unit based}:  
tCO$_2$ per tonne of steel, cement, aluminium, etc.

\paragraph{(3) Value-based}:  
tCO$_2$ per \$ of imported product.

Recommendation:  
Hybrid approach:
\begin{itemize}
    \item category or physical-unit for EITE sectors,
    \item value-based for general goods.
\end{itemize}

%%%%%%%%%%%%%%%%%%%%%%%%%%%%%%%%%%%%%%%%%%%%%%%%%%%%%%%%%%%%%%%%%%%%%%%%%%%%
\subsection{C.4 Export Rebates}

Exports are zero-rated.  
Producers still receive the feebate but pay no consumption charge.

Optional variants:
\begin{itemize}
    \item input-level rebates (VAT-style),
    \item benchmark-based rebates,
    \item sector-level fixed rebates.
\end{itemize}

Recommended:  
\textbf{Pure zero-rating + feebate retention} (simplest, WTO-compatible).

%%%%%%%%%%%%%%%%%%%%%%%%%%%%%%%%%%%%%%%%%%%%%%%%%%%%%%%%%%%%%%%%%%%%%%%%%%%%
\subsection{C.5 Administrative Steps}

\begin{enumerate}
    \item Importer submits emissions data (verified or default).
    \item Customs applies the consumption charge.
    \item Domestic producers apply feebate and charge consumption tax.
    \item Exporters zero-rate the consumption charge.
    \item Government recycles revenue to households or remediation functions.
\end{enumerate}

%%%%%%%%%%%%%%%%%%%%%%%%%%%%%%%%%%%%%%%%%%%%%%%%%%%%%%%%%%%%%%%%%%%%%%%%%%%%
\section{Appendix D: Notes on WTO Compatibility}

Although not elaborated in the main text, OECP is consistent with:
\begin{itemize}
    \item GATT Article II:2(a): border adjustments for internal taxes,
    \item GATT Article III: national treatment,
    \item The SCM Agreement: feebates are non-discriminatory,
    \item Established VAT jurisprudence (e.g.\ EU CJEU decisions).
\end{itemize}

A full WTO-law analysis is left for future work.

%%%%%%%%%%%%%%%%%%%%%%%%%%%%%%%%%%%%%%%%%%%%%%%%%%%%%%%%%%%%%%%%%%%%%%%%%%%%
% END OF PART III
%%%%%%%%%%%%%%%%%%%%%%%%%%%%%%%%%%%%%%%%%%%%%%%%%%%%%%%%%%%%%%%%%%%%%%%%%%%%
\end{document}